\documentclass[12pt,a4paper]{article}
\usepackage[utf8]{inputenc}
\usepackage[T1]{fontenc}
\usepackage[spanish]{babel}
\usepackage{geometry}
\geometry{margin=2.5cm}

\title{Hito 2: Preparación y gestión de datos}
\author{Plantilla}
\date{Marzo 2026}

\begin{document}
\maketitle

\section{Objetivo}
Construir y validar la arquitectura de datos para el caso de churn.

\section{Estructura de contenido}
\begin{itemize}
  \item Configuración de entorno y control de versiones.
  \item Infraestructura de datos y gobernanza.
  \item Capa bronze (ingesta y auditoría).
  \item Capa silver (calidad, cuarentena, históricos).
  \item Capa gold (features y tabla spine).
  \item Ejecución del pipeline (Triggered/Continuous).
  \item Evidencias y validaciones.
\end{itemize}

\section{Mermaid pendiente}
\begin{verbatim}
flowchart LR
    A[Fuentes] --> B[Bronze]
    B --> C[Silver]
    C --> D[Gold]
    D --> E[Feature Store]
\end{verbatim}

\section{Checklist}
\begin{itemize}
  \item Entorno colaborativo configurado.
  \item Capas bronze/silver/gold documentadas.
  \item Reglas de calidad y cuarentena justificadas.
  \item Evidencias de recuentos y esquemas.
\end{itemize}

\end{document}
